% Transcription — fonds Alexandre Grothendieck, shelfmark 42, batch 1 (pages 1-6).
% Established from the facsimile archives/batches/42/batch-01.pdf (a mirror of
% https://grothendieck.umontpellier.fr/42.pdf), per the skill
% transcribe-grothendieck.
%
% Six pages, all of them mathematics, all in his hand, and the whole folder in
% one batch. Nothing is skipped; there are no versos of other people's paper.
% Grothendieck numbers the sheets 1 to 6 himself, in the bottom-left corner,
% and the archivists' numbering agrees with his.
%
% The folder is one continuous piece. Page 1 defines a "réseau d'anneaux
% locaux" — an ordered set X together with an inductive system of local rings
% (O_x) such that for x <= y the map O_x -> O_y induces an isomorphism
% (O_x)_{p_xy} -> O_y — its morphisms, the categories R of réseaux and S of
% preschemes, and the canonical functor p : S -> R. Page 2 states the theorem
% the whole folder is for: p : S_{0/S} -> R_{/S} is fully faithful on
% preschemes locally of finite type over a locally noetherian S.
%
% Pages 3 to 6 are the proof, and it is a proof about sheaves rather than about
% réseaux. Page 3 states the Proposition — for a sheaf F whose stalk maps
% F_y -> F_x are injective along each closure, the global sections are exactly
% the coherent systems of germs, Gamma(X,F) = lim_<- F~ — and proves Lemme 1,
% that the locus where a section agrees with a given system is open. Page 4
% draws Corollaire 1 for O_X and states Lemme 2, the commutative-algebra fact
% that carries everything: for A noetherian and p prime there is f in A - p
% with A_f -> A_p injective. Page 5 finishes that proof and draws Corollaire 2,
% for the sheaf of germs of S-morphisms X -> Y. Page 6 gives the two-morphism
% variant, Corollaire 3 — an S-morphism is the same thing as a homomorphism
% p(X) -> p(Y) — and closes with a Remarque and the example X = Spec k[T]
% showing the Proposition's condition cannot simply be dropped.
%
% Two readings are worth recording. First, the word governing the Proposition's
% hypothesis, and repeated on pages 3, 5 and 6, is written with its first
% syllable compressed to a single hook; taken alone it could be read
% "injective" or "surjective". It is settled by the use made of it: Lemme 2
% delivers injectivity of A_q -> A_p, and page 5 concludes from it that "F_q ->
% F_p est injectif" verifies "la condition de la prop.". It is transcribed
% "injective" throughout, with a note at the first occurrence. Second, page 2's
% entire justification below "pleinement fidèle" is crossed out with two large
% X's, and the bottom third of page 4 is written into the margin and over
% itself. Both are transcribed as far as they are legible and marked where they
% are not.
%
% The hand is a fast, light pencil throughout, and a good deal of the
% connective prose is left as \ill{} rather than guessed. The mathematical
% statements themselves are legible.
%
% The dating is the inventory's own, for the folder.
%
% Pass: Opus 5 (claude-opus-5), 2026-08-23 — first pass, unchecked against the pages by a human.
\documentclass[11pt,a4paper]{article}
% Le préambule commun à toutes les transcriptions du fonds Grothendieck.
%
% Il tient en une idée : **l'incertitude se compose, elle ne se résout pas.**
% Une transcription qui lisse un mot douteux a détruit la seule chose pour
% laquelle on la faisait — un lecteur doit pouvoir savoir, à chaque endroit,
% ce qui était sur le feuillet et ce qui a été ajouté. Les six macros qui
% suivent sont donc l'appareil critique, et non de la décoration.
%
% Les mêmes commandes sont reconnues par `scripts/render.mjs`, qui produit la
% vue de lecture du site. Une macro ajoutée ici doit l'être aussi là-bas,
% faute de quoi le rendu échouera bruyamment — ce qui est voulu : un rendu
% silencieusement faux est pire qu'un rendu absent.

\ProvidesPackage{grothendieck}[2026/08/08 Transcriptions du fonds Grothendieck]

\usepackage{amsmath,amssymb,amsthm}
\usepackage{mathrsfs}
\usepackage{stmaryrd}
% Les diagrammes commutatifs sont partout dans ce fonds : c'est la langue
% même de la géométrie algébrique de Grothendieck, et une transcription qui
% les rendrait en prose ne transcrirait rien.
\usepackage{tikz-cd}
% Français par défaut — c'est la langue des feuillets — mais l'anglais est
% chargé aussi : la traduction et le résumé sont des documents anglais, et la
% typographie française (espaces avant les deux-points) y serait une faute.
% Ces documents basculent par \selectlanguage{english} en tête de corps.
\usepackage[english,french]{babel}
\usepackage{fontspec}
\usepackage{geometry}
\geometry{a4paper,margin=25mm,marginparwidth=28mm}
\usepackage{marginnote}
\usepackage{xcolor}
% ulem, not soul. `soul`'s \st and \ul are built on a character-by-character
% scan that XeLaTeX's font machinery breaks: the document fails with an
% undefined control sequence rather than degrading. ulem does the same two jobs
% and is Unicode-engine safe. `normalem` stops it hijacking \emph, which the
% transcriptions use for Grothendieck's own emphasis.
\usepackage[normalem]{ulem}
\usepackage{hyperref}
\hypersetup{colorlinks=true,linkcolor=black,urlcolor=black,pdfborder={0 0 0}}

% --- Métadonnées du lot -----------------------------------------------------
% Reprises telles quelles par le script de rendu, qui les affiche en tête de
% la vue de lecture. Elles fixent ce que le fichier prétend couvrir : sans
% elles, une transcription flotte sans point d'attache dans un fonds de seize
% mille feuillets.

\newcommand{\folder}[1]{\gdef\grothFolder{#1}}
\newcommand{\batch}[1]{\gdef\grothBatch{#1}}
\newcommand{\leaves}[2]{\gdef\grothFirst{#1}\gdef\grothLast{#2}}
\newcommand{\foldertitle}[1]{\gdef\grothFoldertitle{#1}}
\gdef\grothFolder{?}\gdef\grothBatch{?}\gdef\grothFirst{?}\gdef\grothLast{?}\gdef\grothFoldertitle{}

% --- Le résumé --------------------------------------------------------------
%
% Il ouvre la lecture modernisée, sous le titre « Résumé », et il est composé
% en petit. Ce n'est pas un ornement : c'est le seul passage du document qui
% ne soit pas des mathématiques, et un lecteur doit pouvoir le reconnaître
% avant de l'avoir lu — puis le sauter s'il connaît déjà le sujet, ou s'y
% arrêter s'il ne le connaît pas. Le filet qui le suit marque la reprise du
% corps.

\newenvironment{resume}
  {\small\setlength{\parskip}{0.45em}}
  {\par\medskip\hrule\medskip}

% --- Mots-clés ---------------------------------------------------------------
%
% En anglais, en clôture du résumé de la lecture modernisée : le vocabulaire
% moderne sous lequel chercher ce que les feuillets construisent. Le manifeste
% du site (`npm run manifest`) les extrait pour étiqueter la cote entière —
% cette ligne est la source unique des tags, il n'y a pas de fichier à côté.
\newcommand{\keywords}[1]{\par\smallskip\noindent
  {\footnotesize\textsc{Keywords} --- \textit{#1}}\par}

% --- L'appareil critique ----------------------------------------------------

% Le numéro de feuillet, en marge. C'est l'ancre qui permet de régler un
% désaccord sur une formule en regardant *un* feuillet plutôt qu'une chemise
% de six cents. Le site s'en sert aussi pour tourner les pages du fac-similé
% quand on fait défiler la transcription.
\newcommand{\leaf}[1]{%
  \marginnote{\footnotesize\color{gray!70}\textsf{#1}}%
  \label{leaf:#1}\ignorespaces}

% Illisible : on ne devine pas. Un mot inventé qui se lit comme les autres est
% le pire résultat possible.
\newcommand{\ill}{\textcolor{red!60!black}{[\,\ldots\,]}}

% Lecture douteuse : le mot est donné, et signalé comme incertain.
\newcommand{\uncertain}[1]{\uline{#1}}

% Ajout de l'éditeur — une lettre manquante, un mot sous-entendu.
\newcommand{\add}[1]{\textcolor{blue!50!black}{[#1]}}

% Biffé par Grothendieck lui-même. Ce qu'il a rayé fait partie du document :
% une rature dit souvent où la pensée a changé de direction.
\newcommand{\struck}[1]{\sout{#1}}

% Note du transcripteur, sur l'état matériel du feuillet.
\newcommand{\note}[1]{\par\smallskip{\small\color{gray!60!black}\textit{[#1]}}\par\smallskip}

% Note marginale de Grothendieck, distincte du corps du texte.
\newcommand{\marginal}[1]{\par\smallskip\noindent\hspace{1em}%
  {\small\color{gray!60!black}#1}\par\smallskip}

% --- Titre ------------------------------------------------------------------

\AtBeginDocument{%
  \begin{center}
    {\large\bfseries Fonds Alexandre Grothendieck --- cote n\textsuperscript{o}~\grothFolder}\\[2pt]
    {\itshape\grothFoldertitle}\\[4pt]
    {\small feuillets \grothFirst--\grothLast\ (lot \grothBatch)}\\[2pt]
    {\footnotesize\color{gray!60!black}Transcription établie d'après le fac-similé numérisé
     par l'Université de Montpellier.\\
     Les lectures incertaines sont soulignées, les illisibles notées [\,\ldots\,],
     les ajouts entre crochets.}
  \end{center}
  \vspace{1em}}

\endinput


\folder{42}
\batch{1}
\pages{1}{6}
\dating{[années 1960-1970]}
\watermark{Édition de démonstration}
\foldertitle{Réseaux d'anneaux locaux : notes manuscrites (s.d.)}

\begin{document}

\section*{« Réseau d'anneaux locaux » (pages 1 à 6)}

\page{1}
\emph{Réseau d'anneaux locaux.}

\ill{} est formé d'un ensemble ordonné $X$, et \struck{pour tout $x \in X$,}
et d'un système inductif d'anneaux locaux $(O_x)_{x \in X}$, [donc si
$x \leq y$, \ill{} $O_x \to O_y$) tel que pour $x \leq y$, $O_x \to O_y$
induise un isomorphisme

\[ (O_x)_{p_{xy}} \longrightarrow O_y, \]

(où $p_{xy}$ est l'image \ill{} de $\mathfrak{m}(O_y)$))
\struck{et que $p_{xy} \supset p_{xy'} \Rightarrow y \geq y'$.}

\marginal{On obtient ainsi une correspondance \ill{} \uncertain{bijective}
entre les $y \geq x$ et les idéaux premiers de $O_x$.}

Notation $X$, $Y$ \ill{} base du réseau.

Un homom.\ d'un réseau $X$ dans un autre $Y$ : \ill{} des systèmes inductifs
d'anneaux \emph{locaux}, i.e.\ formé d'une application \struck{\ill{}}
croissante

\[ X \xrightarrow{\ f\ } Y, \]

et pour tout $x \in X$ d'un homom.\ \emph{local} $O_{f(x)} \to O_x$,
compatible avec les applications de \struck{spécialisation} généralisation.

Soit $\underline{R}$ la catégorie des réseaux, $\underline{S}$ la catégorie
des préschémas. On a un foncteur canonique

\[ \underline{S} \xrightarrow{\ p\ } \underline{R}. \]

Réseau sur un anneau, sur un préschéma.

Soit $S$ un préschéma loc.\ noethérien,

\page{2}
soient $\underline{S}_{/S}$ et $\underline{R}_{/S}$ les catégories des
préschémas (resp.\ réseaux) sur $S$, $\underline{S}_{0/S}$ la sous-catégorie
de $\underline{S}_{/S}$ formée des préschémas loc.\ de type fini. Alors
\struck{$p$}

\[ p : \underline{S}_{0/S} \longrightarrow \underline{R}_{/S} \]

est \emph{pleinement fidèle}.

\note{Tout ce qui suit sur cette page — la justification, depuis « En effet »
jusqu'à « $V$ de $y$ » — est barré de deux grandes croix. La démonstration
reprend à neuf page 3, par la proposition sur les faisceaux. Ce qui reste
lisible sous les croix est transcrit ci-dessous.}

\struck{[En effet, un homom.\ de réseaux $p(X) \to p(Y)$ est
\uncertain{certainement} \ill{} pour les applications de \ill{} \ill{}, car
\ill{} \ill{} on peut, ici $X$ est affine donc noethérien, \ill{} les pts de
$X$ par spécialisation, d'un ens.\ fini \ill{}, donc \ill{} en pts de $=$ les
pts de $f(X)$ \ill{}}

\struck{Plus généralement, si $X$ est loc.\ de type fini sur $S$,
\ill{}, alors}

\struck{$\mathrm{Hom}_S(X, Y) \xrightarrow{\ \sim\ }
\mathrm{H}_S(p(X), p(Y))$}

\struck{[bijectif]. En effet, soit $x \in X$, $y = f(x)$, un homom.\ local
$O_y \to O_x$ ; il définit un homom.\ d'un voisinage $U$ de $x$ dans un
voisinage $V$ de $y$.}

\page{3}
\emph{Proposition.} Soit $X$ un espace loc.\ noeth.\ dont toute partie fermée
irréductible a exactement un point générique. \struck{Soit} $F$ un faisceau
sur $X$, satisfaisant la condition suivante : pour tt $x \in X$, il existe un
\struck{\ill{}} voisinage \struck{\ill{}} $U$ de $\bar x$ tel que pour
tt $y \in \bar x \cap U$, l'application naturelle $F_y \to F_x$ soit
\emph{injective}.

\note{Le mot est écrit d'un seul trait, la première syllabe réduite à un
crochet ; isolé, il se lirait aussi bien « surjective ». C'est l'usage qui en
est fait qui tranche : le lemme 2 de la page 4 fournit l'injectivité de
$A_q \to A_p$, et la page 5 en conclut que « $F_q \to F_p$ est injectif »
vérifie « la condition de la prop. ». Il est donc lu « injective » ici et
partout ailleurs.}

\struck{Alors} \struck{Soit} Ordonnons $X$ par $y \leq x \Leftrightarrow
\bar y \ni x$, soit $\widetilde{F}$ le \struck{faisceau} système
\struck{inductif} projectif sur $X$ défini par $F$, considérons
l'application naturelle

\[ \Gamma(X, F) \longrightarrow \varprojlim \widetilde{F} \]

Cette application est \emph{bijective}.

La question est locale. On peut donc supposer $X$ noethérien.
Soit \struck{\ill{}} \ill{} pour \ill{} $\to$ \ill{} est surjective, soit
$(\xi_x)_{x \in X} \in \varprojlim F$.

\emph{Lemme 1.} Soit $s$ une section de $F$ sur $U$. Alors l'ens.\ des pts
$x \in U$ tels que $s_x = \xi_x$ est \emph{ouvert}.

S'il contient $x$, il contient les $y \leq x$ i.e.\ les \ill{}
$\bar y \ni x$ ; il suffit de montrer qu'il est \emph{constructible}, donc que
s'il contient $x$, il contient \uncertain{les pts d'un ouvert non vide} de
$\bar x$. Or il existe un \uncertain{ouvert non vide} $U$ de $\bar x$ tel que
pour $y \in U \Rightarrow F_y \to F_x$ injectif, d'où $s_x = \xi_y$. O.K.

La démonstration est \struck{\ill{}} \ill{} \uncertain{finie}.

\page{4}
\emph{Corollaire 1.} Soit $X$ un préschéma loc.\ noeth. Alors
$\Gamma(X, \mathcal{O}_X)$ s'identifie à l'ens.\ des
$(\xi_x)_{x \in X} \in \prod_{x \in X} O_x$ tels que $x \in \bar y$ implique
que $\xi_y$ \uncertain{soit} l'image par généralisation de \struck{$x$}
$\xi_x$.

Il reste à vérifier que les conditions de la prop.\ sont vérifiées, i.e.\
pour le

\emph{Lemme 2.} Soient $A$ un anneau noethérien, $p$ un idéal premier de $A$.
Alors il existe $f \in A - p$ tel que $A_f \to A_p$ soit \emph{injectif} ; si
\struck{\ill{}} dans $q$ est un idéal premier de $A_f$ [i.e.
\struck{provenant} d'un idéal $q' \supset p$ ne rencontrant pas $f$],
\struck{$\Rightarrow$} l'homom.\ $(A_f)_q \to A_p$ est aussi injectif.

\marginal{\uncertain{Modulo}}

\emph{Démonstration du lemme 2.} Soit $J$ le noyau de $A \to A_p$
[i.e.\ \uncertain{l'annulateur} de $A - p$], il \struck{\ill{}} \ill{}
fini de générateurs $x_i$, donc il existe un $f \in A - p$ annulant $J$,
\ill{} \struck{et} alors, $A_f \to A_p$ est injectif.

\note{Le bas de la page est écrit dans la marge et par-dessus lui-même, les
lignes serrées les unes dans les autres ; il porte la vérification de cette
injectivité. Seuls des fragments s'en laissent lire, et ils sont donnés
ci-dessous dans l'ordre où ils viennent.}

En effet, soit \ill{} $x/f^n$ \ill{} $A_p$, $x \in A - p$, \ill{} $n \geq 0$,
i.e.\ \ill{} $A_p$, donc $x$ \ill{} $A_f - p A_f$, \ill{} $y/f^m$ \ill{}
$A_f$, $y \in A - p$ \ill{} $y x f^N = 0$ pour $N$ grand, \ill{}

\page{5}
\struck{\ill{} soit $x/s \in A$ \ill{} $(A_f)_q$ \ill{}}

Pour la dernière assertion, on peut supposer $f = 1$ \struck{\ill{}} i.e.\
$A_f = A$ \struck{\ill{}} [car $A_p = (A_f)_{p A_f}$], i.e.\ on \ill{}
suppose $A \to A_p$ injectif.

\uncertain{Prouvons} $A_q \to A_p$ injectif si $q \supset p$. Soit $x/s$
($x \in A$, $s \in A - q$) dans le noyau, \struck{il reste} \ill{}
\struck{\ill{} on est \ill{} dans $A_p$, et \ill{} injectivité \ill{}
\ill{} $A_p$} il s'ensuit que $x$ est nul, donc $x/s$ est nul \ill{} dans
$A_q$.

\emph{Corollaire 2.} Soit $X$ un préschéma loc.\ noethérien,
\struck{de type fini} sur $S$, soit $Y$ un $S$-préschéma, soit $F$ le
\struck{\ill{}} faisceau sur $X$ des germes de $S$-homomorphismes $X \to Y$.
Alors la donnée d'un $S$-\uncertain{morphisme} $X \to Y$ équivaut à la donnée
d'un système cohérent d'éléments

\[ f_x \in F_x \qquad (x \in X). \]

Il suffit de prouver le

\emph{Lemme 3.} Avec les notations précédentes, le faisceau $F$ satisfait la
condition de la prop.

\marginal{Pour ceci, on peut \ill{} $S$.}

\struck{De} Plus précisément, si $X = \mathrm{Spec}(A)$, $x = p$, et si
$A \to A_p$ est injectif, alors pour les $q$ \uncertain{générisant}
\struck{$x$} $p$, $F_q \to F_p$ est injectif.

\page{6}
\struck{\ill{}} On peut \struck{\ill{}} énoncer : pour \ill{} le résultat un
peu (compliqué) $A$ \ill{} l'\ill{} $A_q$ d'un voisinage de $y$, pour lequel
la condition d'injectivité \ill{} soit satisfaite) :

si $f, g : \mathrm{Spec}(A) \longrightarrow Y$ sont deux morphismes,
\struck{qui} dont les composés avec $\mathrm{Spec}(A_p) \to \mathrm{Spec}(A)$
sont identiques, ils sont identiques.

\emph{Corollaire 3.} Supposons que $X$ soit loc.\ noeth., $Y$
\struck{\ill{}} de type fini sur $S$, $S$ loc.\ noeth.\ (ce dernier
\struck{est}-il nécessaire ?). La donnée d'un $S$-morphisme
\struck{$X \to Y$} équivaut alors à la donnée d'un \ill{} de $p(X)$ dans
$p(Y)$.

\emph{Remarque.} Dans la proposition (et le cor.\ 1 seulement), la
condition \ill{} : \uncertain{sinon} \ill{} pas surj.\ de \ill{}

Ex.\ $X = \mathrm{Spec}\, k[T]$ ($k$ alg.\ clos), \struck{toute} \struck{\ill{}}
\ill{} les pts fermés $x = a$ ($a \in k$), \ill{} on considère
\struck{\ill{}}

\[ F = \coprod_{x \in X} E_x. \]

\ill{} les $F_\varphi$ \ill{} ($\varphi$ \ill{} généralisation \ill{}). Soit
$\xi_x \in \Gamma(E_x)$ l'élément « unité » de $\Gamma(E_x)$, et $\xi_a = 0$ ;
les $\xi_x$ ($x \neq x_i$) forment un syst.\ cohérent d'éléments \ill{}

\end{document}
